She had stopped waiting on the third month, not because she’d given up hope but because she’d worked out, sitting in a crowded hall in White Harbor with four hundred other people doing the same waiting, that hope wasn’t a thing you could simply hold still and let do its work. It wanted feeding, same as anything else that was meant to keep growing, and sitting in a borrowed room waiting for a letter that plainly wasn’t coming had started to feel less like patience and more like a slow, comfortable way of starving it.

So she paid a fisherman named Otter three days’ worth of mending and half her coin for passage on a boat going south with salted cod, and she did not ask Lord Manderly’s leave to go, because she had learned a long while back, beyond a Wall that no longer meant what it once had, that the people with the authority to grant leave were very often the last people who understood what a body actually needed to do.

Little Sam cried for the first two days of the crossing and slept through the last four, and Gilly spent the whole of it teaching him the names of things on the deck — rope, sail, gull, the grey endless water — because she had decided, somewhere in the first month of the waiting, that a boy ought to have words for the world even in the middle of not knowing what was going to happen to either of them.

\sceneBreak

The city came up out of the haze on the eighth morning, and she stood at the rail with the boy on her hip and looked at it for a long while before she let herself feel anything about what she was looking at.

She had heard the size of it described a hundred times in a hundred mouths at White Harbor, and none of the telling had prepared her for the seeing, though she found, standing there, that she was rather less undone by it than she’d expected to be. She had grown up in a place where the world ended a hundred yards past your own fire and everything past that was cold and hungry and trying to kill you. A burned city, however large, was still just a place that had been hurt, and she had spent her whole life around things that had been hurt, and she found she knew, without having to think hard about it, roughly how to hold herself in front of one.

\textit{It’s only a bigger version of everywhere else,} she thought. \textit{Bigger walls, bigger fires, bigger grief. The bones of it are the same as anyplace I’ve ever been. People needing feeding. Children needing carrying. Somebody, somewhere, needing to be found.}

\sceneBreak

Finding Sam took the better part of a day, which surprised her less than it perhaps should have, given that nobody at the harbor gate seemed to have the first idea who governed this city or where a person went to ask after anyone in it.

“Maester Tarly,” the third steward said, when she’d finally worked her way to someone who recognized the name rather than simply shrugging at it. “He’ll be at the tower, or the small hall, or with the queen, depending on the hour.”

“Which tower.”

“The Hand’s tower, my lady. Two hundred and eleven steps, if you’re of a mind to climb them, though I’d leave the boy with somebody first if I were you — it’s a long way up for small legs.”

She did not leave the boy with anybody. She carried him up every one of the two hundred and eleven steps herself, because she had carried him considerably further than that already, across a whole kingdom of ruin, and she was not about to hand him to a stranger for the last small stretch of it.

\sceneBreak

She found Sam sitting on the floor of a room at the top of the tower, beside a bed with a boy in it who was breathing steadily and looking at nothing at all, and she stood in the doorway a moment before he noticed her, because she found, watching him, that he had changed in ways she had not been prepared for.

He looked older. Not simply tired — she had expected tired, everyone she’d passed in this city looked tired — but older in the particular way a man looks when he has spent a year being needed by considerably more people than any one man ought reasonably to be needed by, and she felt, standing there, an unexpected flare of something close to anger at all of them, at this whole city, for having taken so much of him before she’d had the chance to say a proper goodbye to whatever part of him they’d taken it from.

“Sam.”

He turned, and for a moment his face did not do anything at all, the particular stillness of a man who has imagined a scene so many times that the real arrival of it does not immediately register as real, and then it did register, all at once, and he was on his feet and across the room before she’d finished saying his name a second time.

\sceneBreak

He held her for a long while, and held the boy after, lifting him with the particular careful astonishment of a man discovering that a child he’d last held as an infant had turned, somewhere in the intervening year, into an actual small person with opinions and a considerable quantity of new words.

“You’re enormous,” he told little Sam, who regarded this pronouncement with the grave suspicion small children reserve for adults saying obvious things as though they were remarkable. “You were the size of a loaf of bread. Look at you.”

“He walks now,” Gilly said. “Has done since the spring. Talks a fair bit too, mostly about gulls, on account of that’s near all he saw for the last eight days.”

“I didn’t know that. I didn’t know any of it.” Sam’s voice caught slightly. “I sent word. Twice. I don’t know if either letter ever reached you, and I’ve had no answer to either, and I’ve spent months not knowing whether that meant you’d never had them or—“

“I never had them,” Gilly said. “Whatever you sent, it never reached White Harbor, or if it did, it went the way most things went there this past year, into a pile nobody had the hands to sort.” She said this plainly, without much heat, because she had decided on the crossing south that there was no profit in being angry at a silence neither of them had chosen. “I stopped waiting on the third month. Not because I’d given up on you. Because waiting wasn’t doing either of us any good, and I’ve never yet met a problem that got smaller for being sat next to.”

\sceneBreak

She met the fat one they called Maester Aldrich two days later, and Sansa Stark the day after that, and found she had opinions about both of them that she did not trouble to keep entirely to herself.

“You’ve a whole city keeping records on wax and paper,” she told Sam, on the third evening, watching him work through a stack of the day’s business at the small trestle, “and not one of you thought to keep a proper count of who’d been sent where when the North emptied out. I’d have had that sorted in a week with four women and a ledger, back home, and we’d none of us had your letters or your maesters.”

“You’re not wrong.”

“I know I’m not wrong. I’d tell you so whether you agreed or not.” She said it without malice, in the same flat practical register she said most things. “It’s not a criticism of you particularly, Sam. It’s a criticism of a city that’s spent a year being clever about everything except the one thing that actually would have found me eight months sooner.”

\sceneBreak

She found her way, in the weeks that followed, into the actual working life of the city rather more quickly than she’d expected to, because a woman who had delivered half the babies born beyond a Wall that no longer stood for anything was not a woman this particular winter had any use in wasting.

She helped where help was wanted, in the birthing rooms mostly, because that was the trade she actually knew, and found the maesters and the highborn midwives of King’s Landing did things by the book in a manner she’d rarely had the luxury of doing anything, and found also, within a fortnight, that the book was wrong about a good deal more than anybody south of the Wall seemed prepared to admit.

“You’d have her lie flat the whole while,” she told a King’s Landing midwife, watching a difficult birth go worse than it needed to in a house off the Street of Sisters, “and flat’s the worst position there is for a woman whose hips don’t want to open the way they should. Get her up. Hands and knees, or squatting if she’s the strength for it. Gravity does half the maester’s job for him if you’ll only let it, and I’ve pulled more babies out alive on hands and knees than I ever did with a woman on her back staring at a ceiling.”

The midwife did not entirely believe her, that first time. Gilly did not much mind. She had learned a long while back that a thing proved once was worth considerably more than a thing merely argued, and she went on offering what she knew, plainly, to whoever would take it, and let the results do the arguing she couldn’t be bothered to do herself.

\sceneBreak

She thought, some evenings, climbing the two hundred and eleven steps with Sam to sit awhile with the boy who breathed and did not answer, that she had come an enormous distance to end up doing precisely the same work she’d done her whole life — feeding people, carrying children, being the steady hands in a room when somebody’s steady hands were the only thing standing between a birth going well and a birth going very badly indeed.

“I didn’t come here to be a maester’s wife sitting quiet in a tower,” she told Sam, one evening, plainly, because she had decided a long while back that plainness was the only thing she had ever been able to properly afford. “I came here because you were the one thing left in this world I hadn’t given up on finding. But I’ll not sit idle while I’m here either, Sam. There’s women in this city dying of things I’ve watched other women survive a hundred times over with considerably less than what’s sitting in these walls. I mean to be of use, and not simply of comfort, and I’d thank you not to be surprised by that, given everything you already know about me.”

“I’m not surprised,” Sam said. “I don’t think I’ve ever once been surprised by that about you. I think it’s rather the whole of why I wrote your name in the wax before I’d learned to write anything else properly at all.”

Gilly did not entirely know what to make of that, and did not ask him to explain it, and simply took his hand instead, ink-stained and steady, and held it the way she had held it across a great many worse nights than this one, and found that it still fit exactly the way she remembered.